\section{讨论}
\label{sec:discussion}

在理论上，扩大病例规模能换来多少精度取决于估计目标。瞬时再生数的规模收益不饱和，方差随规模持续下降（命题~\ref{prop:current}）。跨期平均再生数的误差含两类不随规模下降的方差。环境方差 $v_R$ 来自传播条件的逐期波动，报告异质性方差 $v_H$ 来自报告比例的逐期波动（命题~\ref{prop:average} 与命题~\ref{prop:report-noise}）。方差平衡规模给出了规模收益开始明显减小的病例数，报告率波动时它由 $A/v_R$ 降为 $A_H/(v_R+v_H)$。Parag 等\citep{parag2022quantifying, parag2022fundamental}的信息论分析刻画了随规模下降的那一部分，在离散参数趋于无穷时与命题~\ref{prop:current} 一致。生态学中环境随机性与人口随机性的区分\citep{engen1998demographic, rohani2002determinism}由此在监测设计中有了一个以病例数计的对应量，报告异质性则是监测系统自身带来的第二个不可约成分。仅凭计数的前两阶矩无法区分这两类方差，区分它们需要报告率的辅助数据。

上述下确界不约束允许条件偏差的估计量。收缩估计、经验贝叶斯估计\citep{efron1975data}以及 EpiFilter 一类借用相邻时期信息的平滑方法\citep{parag2021improved}，都以对 $R_t$ 的条件偏差换取方差。以补充材料 S1.11 节的线性收缩为例，取图~\ref{fig:variance} 的参数与 $m=50$，$V=0.0825$；先验中心偏离 $\Rbar$ 不超过 0.279 时，最优收缩的均方误差就低于 $v_R=0.04$，偏离 0.1 时为 0.0089，约为 $v_R$ 的 22\%。代价是估计量向先验中心偏移，偏差 $\lambda(\mu_0-R_t)$ 随当期环境变化，在 $R_t$ 远离常态的时期最大，而这正是预警最需要准确的时候。$R_t$ 有时间相关时，平滑方法还能利用 $(1-\phi^2)v_R<v_R$ 这一更低的条件方差。本文关于下确界与方差平衡规模的结论只覆盖条件无偏的单期估计和 H4 下的一步预测，不适用于这些估计量；它们的均方误差随规模如何变化，需要另行推导。

在监测设计上，三类杠杆对应不同的投入。扩大单期规模降低抽样项，直到误差接近 $v_R+v_H$ 构成的平台；它对预警始终有价值，对趋势估计的价值在规模达到方差平衡规模的量级后明显减小。增加观测期数或汇总更多辖区降低环境项；规模足够大以后，$v_R$ 本身的估计精度也只随期数提高（命题~\ref{prop:vR-info}）。统一报告口径，或在估计中使用观测到的报告率，去除 $v_H$。在命题~\ref{prop:levers} 的数值例子中，增加期数的效果大于同等倍数地扩大规模（图~\ref{fig:r3-levers}a）。报告偏差在各期持续时，增加期数对报告项无效，只有统一口径能处理。期数与规模之间的取舍还取决于环境的时间相关。$R_t$ 正相关时，预算低于阈值 $\mathcal B^\dagger$ 应设更少、更大的期，高于阈值才应设更多、更小的期（命题~\ref{prop:ar-design}）。独立情形下 $m^*$ 的近似只在大预算下成立，按它设计小预算监测会选出过小的单期规模。

对监测实践，预警和趋势估计宜采用不同的设计。判断本周疫情是否在增长属于估计目标1，规模的价值持续存在，小规模时应使用精确检验，并同时报告首次达到和稳定保持目标功效所需的规模（第~\ref{sec:res-test}~节），资源应投向报告及时、覆盖完整的单元。估计一段时期的平均传播水平属于估计目标2，规模到达方差平衡规模的量级以后，增加观测期数或汇总多个辖区更为有效（式~\eqref{eq:ar-sum} 与式~\eqref{eq:pooling}），全国加总序列的结果与空间方向的推论一致（第~\ref{sec:res-robust}~节）。报告率固定且很低时，提高报告率与扩大覆盖的效果相差无几（注~\ref{rem:reporting}），取舍主要取决于成本。报告率随期波动时，二项型波动在规模足够大时使提高报告率占优，制度型波动下扩大覆盖仍占优，两条途径都不改变 $v_H$（推论~\ref{cor:switch}）。NHSN 的填报医院比例在流行初期波动很小，自愿填报期则大了两个数量级（第~\ref{sec:res-reporting-adj}~节）。报告标准化的收益因此主要出现在填报制度变化的时期，切换制度时应同时记录各单位的填报状态，以便事后按报告率调整。

流行阶段转换时出现了另一种误差来源。达峰期以初期均值为基准时，绝大部分损失来自基准偏移，扩大规模无法消除它；换用峰期自身的基准后，规模项重新为正（第~\ref{sec:res-peak}~节）。阶段转换时应优先更新基准模型，同时继续维持病例规模。FluSight 集成预测在峰期的偏差占比较低，提示能跟踪阶段变化的预测模型可以压低这一项；该对比来自不同的模型和样本，只有方向上的参考价值。

报告制度切换的证据需要按流行季解读。全样本上第二期稳健离散度不变的表象，来自第二期季内州周所占比例很低；分层后，季外三期几乎相同，季内第二期的离散度明显扩大（第~\ref{sec:res-policy}~节）。季内证据全部来自自愿填报期的最后一个月，区间很宽，填报不全造成的跳变与该月本身的流行特征无法分开。对监测部门而言，填报制度变化的影响可能集中在病例开始增多的时期，数据质量监控应按流行季分别进行。

在实证解读上，主要局限在于增长统计量的口径。流感的世代间隔短于一周，$G_t$ 是两代左右增长的合成，其均值与方差都不同于单代模型中的 $R_t$。实证分析因此只检验了误差随规模变化的形状，拟合的 $a$ 和 $b$ 不能解读为 $v_R/\Rbar^2$ 与 $A/\Rbar^2$ 的估计，经验比值 $b/a$ 与情景区间的吻合也只在量级上成立；估计单代参数需要日度数据或显式的世代间隔反卷积。均值形状的比较支持误差中含有不随规模消失的部分，但在样本覆盖的规模范围内，严格的底板与指数约 0.27 的幂律衰减拟合相近，更大的规模范围才能区分两者。

截距与斜率的绝对水平都只得到部分解释。斜率比与式~\eqref{eq:b-delta} 的预测吻合，纯测量误差模型给出几乎相同的预测，斜率比因此只表明分母含有报告计数噪声，不能用来判别环境方差；斜率的绝对水平约有一半能由二项分母噪声解释。截距比区间排除了平稳环境模型，但与增长期 $\phi=0$ 的日分辨率更新过程和增长期乘性报告噪声模型的范围相交。入院延迟在环境各期独立时也会产生 $q_1<1<q_2$，因此截距比的方向不能单独诊断 $R_t$ 的时间相关；跨病原分级为复现 0、部分复现 4、方向一致 3，证据主要来自截距超出量。在情景范围的环境方差下，时间相关与归一化效应合计约占实证截距的六成（第~\ref{sec:res-inflation}~节）。把世代块从 7 天缩短到 4 天时截距比几乎不变，截距不升反降（补充材料 S2.4 节），一周跨两代不是剩余部分的来源。余下的候选来源是截距的季层面成分与同填报比例无关的报告离散。各季截距之间的差异远大于季内抽样误差（图~\ref{fig:seasons}b），前者的大小留待更多流行季确定；后者会抬高截距而不改变 $\Delta a$，目前没有直接检验。

流行季的数目是本文最根本的限制。按推论~\ref{cor:season}，季层面方差的估计精度由季数决定，增加州数、周数或报告规模都不改变它。实证结果与此一致，截距的重抽标准差在规模外推到无穷大时几乎不变，把季作为随机效应时区间很宽（第~\ref{sec:res-seasons}~节）。正文中按州聚类的区间因此只描述给定已观测三个季时的不确定性。对截距水平的进一步认识需要积累更多流行季，扩大单季样本的作用有限。把第 4 季纳入季随机效应后，相对标准误的下界仍然很高（补充材料 S2.3 节）。跨病原与样本外季的复制都来自同一个 NHSN 文件，共用同一套报告系统，系统层面的报告特征无法由这些复制排除。病例定义不同的 ILINet 门诊监测给出同向的截距比，FluSurv-NET 的站点过少，区间过宽（补充材料 S2.5 节）。

方差分解依赖假设 H5，即报告率 $p_t$ 与再生数 $R_t$ 相互独立。$v_H$ 与 $v_R$ 能分开写出，前提是两者的波动互不相关。若检测或填报行为随流行强度变化，例如高峰期检测更积极或填报更滞后，报告波动会与 $R_t$ 相关，协方差项会混入截距，$v_H$ 与 $v_R$ 的分解不再成立。本文的报告调整只覆盖可观测的填报医院比例，不覆盖随流行强度变化的就诊、检测与收治行为，可观测报告成分贡献很小的结论不能推广到报告行为整体。

推断方法上，区间来自按州聚类的 2,000 次重抽，按州 $\times$ 流行季聚类的结论相同，但本文没有做野聚类重抽。州与周的双向聚类在聚类数较少时方差估计不稳定\citep{cameron2011robust, cameron2015practitioner}，而每个流行初期窗口只有八到十一周。工具变量回归的排除限制无法检验，其结果只能作为替代测度下的稳健性证据。

部分结论依赖外部参数和事后划定的窗口。情景区间依赖离散参数 $k$ 和住院比例 $\rho$ 的外部取值，二者在流感住院数据中都没有精确估计；精确增长检验要求 $k$ 已知，以估计值代入后水平不再有精确保证。由全国序列推算的州间相关忽略了分母噪声和口径问题。事前基准解决了基准的前视问题，窗口选择本身仍是事后的；报告调整、等价检验、规模与期数比较在运行前写定了判据，季随机效应与规模外推则是事后分析。缺失州周按零处理，而未填报与零住院并不等同；剔除含缺失周的州周后，第二期季内的结论不变。

综合来看，病例规模的边际价值由估计目标与两类不可约方差共同决定。预警需要尽可能大的规模；以条件无偏估计或一步预测为目标的趋势估计，在规模达到方差平衡规模的量级后，应把新增资源转向更多观测期、更广的空间汇总、统一的报告口径和及时更新的基准模型。环境方差本身的大小只能靠积累观测期来认识，在本文数据中，这意味着积累流行季。允许偏差的平滑估计可以进一步压低趋势估计的误差，代价是对异常时期的反应变慢。
